%!LW recipe=pdflatex x3
\documentclass[10pt,a4paper]{article}

\usepackage[spanish,es-nodecimaldot]{babel}
\usepackage[utf8]{inputenc}
\usepackage[T1]{fontenc}
\usepackage{newtxtext,newtxmath}
\usepackage{geometry}
\usepackage{microtype}
\usepackage{amsmath,mathtools}
\usepackage{enumitem}
\usepackage{xcolor}
\usepackage{hyperref}
\usepackage{fancyhdr}
\usepackage{lastpage}
\usepackage[most]{tcolorbox}

\geometry{margin=2.5cm,includeheadfoot}
\setlength{\headheight}{15pt}
\setlength{\headsep}{0.28cm}
\setlength{\footskip}{0.62cm}
\setlength{\parindent}{0pt}
\setlength{\parskip}{0.28em}
\setlength{\abovedisplayskip}{0.5em}
\setlength{\belowdisplayskip}{0.5em}
\setlength{\abovedisplayshortskip}{0.25em}
\setlength{\belowdisplayshortskip}{0.25em}
\setlist[enumerate]{leftmargin=2.2em,itemsep=0.3em,topsep=0.25em}
\setlist[itemize]{leftmargin=1.5em,itemsep=0.2em,topsep=0.2em}

\definecolor{repblue}{gray}{0}
\definecolor{repteal}{gray}{0}
\definecolor{repgray}{gray}{0}
\definecolor{repaccent}{gray}{0}
\definecolor{repline}{gray}{0}
\definecolor{replight}{gray}{1}

\hypersetup{
  hidelinks,
  pdftitle={TRGF Padilla Tarea 5},
  pdfauthor={Juan Ignacio Padilla Barrientos}
}

\pagestyle{fancy}
\fancyhf{}
\fancyhead[L]{\textcolor{repblue}{\small\itshape TRGF Padilla}}
\fancyhead[R]{\textcolor{repgray}{\small Seminario UCR I 2026}}
\fancyfoot[C]{\textcolor{repgray}{\small \thepage/\pageref*{LastPage}}}
\renewcommand{\headrulewidth}{0.3pt}
\renewcommand{\footrulewidth}{0pt}

\newcommand{\CC}{\mathbb{C}}
\newcommand{\Hom}{\operatorname{Hom}}
\newcommand{\Id}{\operatorname{Id}}

\tcbset{
  enhanced,
  sharp corners,
  boxrule=0.45pt,
  arc=0pt,
  left=0.85em,
  right=0.85em,
  top=0.45em,
  bottom=0.4em,
  boxsep=0pt
}

\newcommand{\inlinehead}[2][repblue]{%
  \par\smallskip
  {\color{#1}\bfseries #2.}\hspace{0.45em}%
}

\newcommand{\answer}{\inlinehead{Respuesta}}

\newcommand{\topbanner}{%
  \begin{tcolorbox}[
    colback=replight,
    colframe=repline,
    borderline west={1.2pt}{0pt}{repblue},
    borderline east={1.2pt}{0pt}{repteal},
    left=1.0em,
    right=1.0em,
    top=0.7em,
    bottom=0.55em
  ]
    {\Large\bfseries\color{repblue} TRGF Padilla \textemdash\ Tarea 5}\hfill
    {\color{repaccent}\small\itshape Soluciones}\\[-0.1em]
    {\normalsize\color{repteal} Seminario UCR I 2026}
  \end{tcolorbox}
}

\newcounter{probcnt}
\newcount\workspacelines

\newcommand{\problem}[2]{%
  \stepcounter{probcnt}%
  \ifnum\value{probcnt}>1\newpage\fi
  \par\vspace{0.25em}
  {\large\bfseries\color{repblue} Problema #1}\par
  {\color{repline}\rule{\linewidth}{0.55pt}}\par
  \begin{tcolorbox}[
    breakable,
    colback=white,
    colframe=repline,
    borderline west={1.6pt}{0pt}{repteal},
    left=0.7em,
    right=0.7em,
    top=0.35em,
    bottom=0.3em,
    before skip=0.35em,
    after skip=0.45em
  ]
    #2
  \end{tcolorbox}
}

\newcommand{\workspace}[1]{%
  \answer
  \par\vspace{0.35em}
  \begingroup
  \workspacelines=0
  \loop
    \ifnum\workspacelines<#1
      \vspace{0.72cm}
      \advance\workspacelines by 1
  \repeat
  \endgroup
  \vspace{0.35em}
}

\begin{document}

\color{black}
\topbanner
\small

\problem{1}{%
(*) Demuestre que la matriz del producto tensorial de dos matrices es lo que dijimos en clase.
}

\answer
\medskip

\textbf{Definición (producto de Kronecker).} Dadas matrices $A=(a_{ki})\in M_{m\times n}(k)$
y $B\in M_{p\times q}(k)$, su \emph{producto de Kronecker} es la matriz de bloques
\[
  A\otimes B
  :=
  \begin{pmatrix}
    a_{11}B & a_{12}B & \cdots & a_{1n}B \\
    a_{21}B & a_{22}B & \cdots & a_{2n}B \\
    \vdots  &         & \ddots & \vdots  \\
    a_{m1}B & a_{m2}B & \cdots & a_{mn}B
  \end{pmatrix}
  \in M_{mp\times nq}(k).
\]

\medskip
\textbf{Demostración.}
Sean $(U,\pi)$ y $(V,\rho)$ representaciones de $G$ sobre $k$, con bases
$\mathcal{B}=\{e_1,\ldots,e_m\}$ y $\mathcal{C}=\{f_1,\ldots,f_n\}$ respectivamente.
El espacio $U\otimes V$ tiene base
\[
  \mathcal{B}\otimes\mathcal{C}
  =\{e_i\otimes f_j : 1\le i\le m,\;1\le j\le n\},
\]
ordenada lexicográficamente:
$(e_1\otimes f_1,\,e_1\otimes f_2,\ldots,e_1\otimes f_n,\;e_2\otimes f_1,\ldots,e_m\otimes f_n)$.

Fijemos $g\in G$ y sean $A=[\pi(g)]_{\mathcal{B}}$ y $B=[\rho(g)]_{\mathcal{C}}$,
con entradas $a_{ki}$ y $b_{lj}$.
La representación producto tensorial actúa por
\[
  (\pi\otimes\rho)(g)(e_i\otimes f_j)
  =\pi(g)(e_i)\otimes\rho(g)(f_j).
\]
Escribiendo explícitamente en bases,
\[
  \pi(g)(e_i)=\sum_{k=1}^m a_{ki}\,e_k,
  \qquad
  \rho(g)(f_j)=\sum_{l=1}^n b_{lj}\,f_l,
\]
se obtiene
\[
  (\pi\otimes\rho)(g)(e_i\otimes f_j)
  =\sum_{k=1}^{m}\sum_{l=1}^{n} a_{ki}\,b_{lj}\;(e_k\otimes f_l).
\]
Por lo tanto, la entrada en la fila $(k,l)$ y columna $(i,j)$ de la matriz
$[(\pi\otimes\rho)(g)]_{\mathcal{B}\otimes\mathcal{C}}$ es $a_{ki}\,b_{lj}$.

Por otra parte, en $A\otimes B$ el bloque $(k,i)$ es $a_{ki}B$, cuya entrada
en la fila $l$ y columna $j$ es también $a_{ki}\,b_{lj}$.
Ambas matrices tienen las mismas entradas, luego
\[
  [(\pi\otimes\rho)(g)]_{\mathcal{B}\otimes\mathcal{C}} = A\otimes B.
\]
\hfill$\blacksquare$

\problem{2}{%
Sean $U$, $V$ y $W$ espacios vectoriales sobre un cuerpo $k$. Demuestre que
\[
  \Hom_k(U \otimes V, W) \cong \Hom_k(U, \Hom_k(V, W)).
\]
}

\answer
Definimos primero una aplicación
\[
  \Phi\colon \Hom_k(U\otimes V,W) \to \Hom_k(U,\Hom_k(V,W))
\]
por la fórmula
\[
  \Phi(T)(u)(v):=T(u\otimes v).
\]
Esta definición tiene sentido porque, fijo $u\in U$, la aplicación
$v\mapsto T(u\otimes v)$ es lineal en $v$, ya que $T$ es lineal. Además,
$u\mapsto \Phi(T)(u)$ también es lineal: si $u_1,u_2\in U$ y
$\alpha,\beta\in k$, entonces para todo $v\in V$,
\begin{align*}
  \Phi(T)(\alpha u_1+\beta u_2)(v)
  &=T\bigl((\alpha u_1+\beta u_2)\otimes v\bigr) \\
  &=T\bigl(\alpha(u_1\otimes v)+\beta(u_2\otimes v)\bigr) \\
  &=\alpha T(u_1\otimes v)+\beta T(u_2\otimes v) \\
  &=\alpha\Phi(T)(u_1)(v)+\beta\Phi(T)(u_2)(v).
\end{align*}
Por tanto, $\Phi(T)\in \Hom_k(U,\Hom_k(V,W))$.

Recíprocamente, sea
\[
  F\in \Hom_k(U,\Hom_k(V,W)).
\]
Definimos una aplicación
\[
  b_F\colon U\times V\to W,
  \qquad
  b_F(u,v):=F(u)(v).
\]
Como $F$ es lineal y cada $F(u)$ es lineal en $V$, la aplicación $b_F$ es
bilineal. En efecto,
\begin{align*}
  b_F(\alpha u_1+\beta u_2,v)
  &=F(\alpha u_1+\beta u_2)(v) \\
  &=(\alpha F(u_1)+\beta F(u_2))(v) \\
  &=\alpha F(u_1)(v)+\beta F(u_2)(v) \\
  &=\alpha b_F(u_1,v)+\beta b_F(u_2,v),
\end{align*}
y, para $v_1,v_2\in V$ y $\lambda,\mu\in k$,
\begin{align*}
  b_F(u,\lambda v_1+\mu v_2)
  &=F(u)(\lambda v_1+\mu v_2) \\
  &=\lambda F(u)(v_1)+\mu F(u)(v_2) \\
  &=\lambda b_F(u,v_1)+\mu b_F(u,v_2).
\end{align*}

Por la propiedad universal del producto tensorial, existe una única aplicación
lineal
\[
  \widetilde{F}\colon U\otimes V\to W
\]
tal que
\[
  \widetilde{F}(u\otimes v)=b_F(u,v)=F(u)(v)
  \qquad
  \text{para todo }u\in U,\;v\in V.
\]
Esto define una segunda aplicación
\[
  \Psi\colon \Hom_k(U,\Hom_k(V,W))\to \Hom_k(U\otimes V,W),
  \qquad
  \Psi(F):=\widetilde{F}.
\]

Veamos que $\Phi$ y $\Psi$ son inversas. Si $T\in \Hom_k(U\otimes V,W)$,
entonces para todo $u\in U$ y $v\in V$,
\[
  \Psi(\Phi(T))(u\otimes v)=\Phi(T)(u)(v)=T(u\otimes v).
\]
Como ambas aplicaciones son lineales y coinciden en todos los tensores puros,
se sigue que $\Psi(\Phi(T))=T$.

Análogamente, si $F\in \Hom_k(U,\Hom_k(V,W))$, entonces para todo
$u\in U$ y $v\in V$,
\[
  \Phi(\Psi(F))(u)(v)=\Psi(F)(u\otimes v)=F(u)(v).
\]
Por consiguiente, $\Phi(\Psi(F))=F$.

Concluimos que $\Phi$ es un isomorfismo, es decir,
\[
  \Hom_k(U \otimes V, W) \cong \Hom_k(U, \Hom_k(V, W)).
\]
\hfill$\blacksquare$

\problem{3}{%
Calcule las matrices que representan a los elementos de $S_3$ en la representación
$V_{\mathrm{st}} \otimes V_{\mathrm{st}}$ en alguna base.

\smallskip
\textit{Para la exposición en clase, me basta con que escoja un elemento de cada clase
conjugada para mostrar.}
}

\answer
Las clases conjugadas de $S_3$ son
\[
  \{e\},
  \qquad
  \{(12),(13),(23)\},
  \qquad
  \{(123),(132)\}.
\]
Por lo tanto, para la exposición basta tomar como representantes
$e$, $(12)$ y $(123)$.

Tomemos una base $(v_1,v_2)$ de $V_{\mathrm{st}}$ en la cual
\[
  \rho((12))=
  \begin{pmatrix}
    -1 & -1 \\
    0 & 1
  \end{pmatrix},
  \qquad
  \rho((123))=
  \begin{pmatrix}
    -1 & -1 \\
    1 & 0
  \end{pmatrix}.
\]
Sea ahora la base de $V_{\mathrm{st}}\otimes V_{\mathrm{st}}$
\[
  \mathcal{B}:=(v_1\otimes v_1,\; v_1\otimes v_2,\; v_2\otimes v_1,\; v_2\otimes v_2).
\]
Por el Problema~1, para todo $g\in S_3$ se tiene
\[
  [\rho\otimes\rho](g)_{\mathcal{B}} = [\rho(g)]\otimes [\rho(g)].
\]

\medskip
\textbf{Representantes de las clases conjugadas.}

Para la identidad,
\[
  (\rho\otimes\rho)(e)=I_4.
\]

Para la transposición $(12)$,
\[
  (\rho\otimes\rho)((12))
  =
  \begin{pmatrix}
    -1 & -1 \\
    0 & 1
  \end{pmatrix}
  \otimes
  \begin{pmatrix}
    -1 & -1 \\
    0 & 1
  \end{pmatrix}
  =
  \begin{pmatrix}
    1 & 1 & 1 & 1 \\
    0 & -1 & 0 & -1 \\
    0 & 0 & -1 & -1 \\
    0 & 0 & 0 & 1
  \end{pmatrix}.
\]

Para el $3$-ciclo $(123)$,
\[
  (\rho\otimes\rho)((123))
  =
  \begin{pmatrix}
    -1 & -1 \\
    1 & 0
  \end{pmatrix}
  \otimes
  \begin{pmatrix}
    -1 & -1 \\
    1 & 0
  \end{pmatrix}
  =
  \begin{pmatrix}
    1 & 1 & 1 & 1 \\
    -1 & 0 & -1 & 0 \\
    -1 & -1 & 0 & 0 \\
    1 & 0 & 0 & 0
  \end{pmatrix}.
\]

\medskip
\textbf{Las seis matrices.}

Además,
\[
  \rho((132))=\rho((123))^2=
  \begin{pmatrix}
    0 & 1 \\
    -1 & -1
  \end{pmatrix},
\]
\[
  \rho((23))=\rho((12))\rho((123))=
  \begin{pmatrix}
    0 & 1 \\
    1 & 0
  \end{pmatrix},
  \qquad
  \rho((13))=\rho((123))\rho((12))=
  \begin{pmatrix}
    1 & 0 \\
    -1 & -1
  \end{pmatrix}.
\]
Por consiguiente,
\[
  (\rho\otimes\rho)((23))
  =
  \begin{pmatrix}
    0 & 0 & 0 & 1 \\
    0 & 0 & 1 & 0 \\
    0 & 1 & 0 & 0 \\
    1 & 0 & 0 & 0
  \end{pmatrix},
\]
\[
  (\rho\otimes\rho)((13))
  =
  \begin{pmatrix}
    1 & 0 & 0 & 0 \\
    -1 & -1 & 0 & 0 \\
    -1 & 0 & -1 & 0 \\
    1 & 1 & 1 & 1
  \end{pmatrix},
\]
\[
  (\rho\otimes\rho)((132))
  =
  \begin{pmatrix}
    0 & 0 & 0 & 1 \\
    0 & 0 & -1 & -1 \\
    0 & -1 & 0 & -1 \\
    1 & 1 & 1 & 1
  \end{pmatrix}.
\]

En resumen, en la base $\mathcal{B}$ ya quedaron calculadas las matrices de
todos los elementos de $S_3$ en $V_{\mathrm{st}}\otimes V_{\mathrm{st}}$, y para la
exposición basta mostrar las de $e$, $(12)$ y $(123)$, que representan las tres
clases conjugadas.
\hfill$\blacksquare$

\problem{4}{%
Encuentre la fórmula general para la descomposición de
\[
  V_{\mathrm{st}}^{\otimes n}
\]
en irreps de $S_3$.
}

\answer
Sean
\[
  \chi_{\mathrm{triv}}=(1,1,1),
  \qquad
  \chi_{\mathrm{sgn}}=(1,-1,1),
  \qquad
  \chi_{\mathrm{st}}=(2,0,-1),
\]
en el orden de clases conjugadas
\[
  \{e\},
  \qquad
  \{(12),(13),(23)\},
  \qquad
  \{(123),(132)\}.
\]
Por Maschke, existen enteros $a_n,b_n,c_n\ge 0$ tales que
\[
  V_{\mathrm{st}}^{\otimes n}
  \cong a_nV_{\mathrm{triv}}\oplus b_nV_{\mathrm{sgn}}\oplus c_nV_{\mathrm{st}}.
\]
Si $W\cong U\oplus V$, entonces la matriz de $g$ en $W$ es por bloques diagonales,
de modo que
\[
  \chi_W(g)=\chi_U(g)+\chi_V(g).
\]
Por consiguiente,
\[
  \chi_n:=\chi_{V_{\mathrm{st}}^{\otimes n}}
  =a_n\chi_{\mathrm{triv}}+b_n\chi_{\mathrm{sgn}}+c_n\chi_{\mathrm{st}}.
\]

Además, si $U$ y $V$ son representaciones, entonces
\[
  (\rho_U\otimes\rho_V)(g)=\rho_U(g)\otimes\rho_V(g),
\]
y como $\operatorname{tr}(A\otimes B)=\operatorname{tr}(A)\operatorname{tr}(B)$,
se sigue que
\[
  \chi_{U\otimes V}(g)=\chi_U(g)\chi_V(g).
\]
Por tanto,
\[
  \chi_n=\chi_{\mathrm{st}}^n=(2^n,0,(-1)^n).
\]

Igualando coordenadas en
\[
  (2^n,0,(-1)^n)
  =a_n(1,1,1)+b_n(1,-1,1)+c_n(2,0,-1),
\]
obtenemos el sistema
\[
  a_n+b_n+2c_n=2^n,
  \qquad
  a_n-b_n=0,
  \qquad
  a_n+b_n-c_n=(-1)^n.
\]
De $a_n-b_n=0$ resulta $a_n=b_n$. Sustituyendo,
\[
  2a_n+2c_n=2^n,
  \qquad
  2a_n-c_n=(-1)^n,
\]
de donde
\[
  a_n=b_n=\frac{2^{n-1}+(-1)^n}{3},
  \qquad
  c_n=\frac{2^n-(-1)^n}{3}.
\]
Así, para todo $n\ge 1$,
\[
  V_{\mathrm{st}}^{\otimes n}
  \cong
  \frac{2^{n-1}+(-1)^n}{3}\,V_{\mathrm{triv}}
  \oplus
  \frac{2^{n-1}+(-1)^n}{3}\,V_{\mathrm{sgn}}
  \oplus
  \frac{2^n-(-1)^n}{3}\,V_{\mathrm{st}}.
\]
Si se adopta la convención usual $V^{\otimes 0}=\CC$, entonces
\[
  V_{\mathrm{st}}^{\otimes 0}\cong V_{\mathrm{triv}}.
\]
\hfill$\blacksquare$

\problem{5}{%
Sea $C$ una clase conjugada y sea
\[
  s_C = \sum_{g\in C} g \in \CC[G].
\]
Sea $(U,\pi)$ una irrep de $G$. Demuestre que
\[
  \pi(s_C) = |C|\,\frac{\chi_U(C)}{\dim U}\,\Id,
\]
donde $\chi_U(C)$ es el carácter de la representación evaluado en cualquier $g$ de esa
clase conjugada.
}

\answer
Sea
\[
  T:=\pi(s_C).
\]
Por el Problema~2 de la Tarea~3, la misma demostración muestra que $s_C$ es
central en $\CC[G]$ para cualquier grupo finito. En consecuencia, $T$ conmuta
con todos los operadores $\pi(h)$, $h\in G$. Entonces, por el mismo argumento
del Problema~2 de la Tarea~4, como $(U,\pi)$ es irreducible se sigue que
$T=\lambda I$ para algún $\lambda\in\CC$.

Solo queda calcular $\lambda$. Tomando trazas en ambos lados,
\[
  \lambda\,\dim U
  =\operatorname{tr}(T)
  =\sum_{g\in C}\operatorname{tr}(\pi(g))
  =\sum_{g\in C}\chi_U(g)
  =|C|\,\chi_U(C),
\]
ya que el carácter es constante en clases conjugadas. En consecuencia,
\[
  \lambda=|C|\,\frac{\chi_U(C)}{\dim U},
\]
y así
\[
  \pi(s_C)=T=|C|\,\frac{\chi_U(C)}{\dim U}\,\Id.
\]
\hfill$\blacksquare$

\problem{6}{%
Sea $\chi_{\mathrm{reg}}$ el carácter de la representación regular, es decir, el álgebra de
grupo pensada como representación de $G$. Demuestre que
\[
  \chi_{\mathrm{reg}}(g) =
  \begin{cases}
    |G|, & \text{si } g = e,\\
    0, & \text{si } g \ne e.
  \end{cases}
\]
}

\answer
Consideremos la representación regular izquierda
\[
  \rho_{\mathrm{reg}}\colon G\to GL(\CC[G]),
  \qquad
  \rho_{\mathrm{reg}}(g)(x)=gx
  \quad (x\in G),
\]
extendida linealmente a $\CC[G]$. Tomemos la base natural de $\CC[G]$
formada por los elementos de $G$.

Para un $g\in G$ fijo, la transformación $\rho_{\mathrm{reg}}(g)$ permuta esta
base, pues envía cada vector de la base $x\in G$ al vector de la base $gx$.
Por tanto, su matriz en esta base es una matriz de permutación y su traza es
el número de vectores de la base que quedan fijos por la multiplicación por $g$.

Si $g=e$, entonces
\[
  \rho_{\mathrm{reg}}(e)=I,
\]
de modo que
\[
  \chi_{\mathrm{reg}}(e)=\operatorname{tr}(I)=\dim \CC[G]=|G|.
\]

Si $g\ne e$, un vector de la base $x\in G$ quedaría fijo solo si
\[
  gx=x,
\]
lo cual implica, multiplicando a la derecha por $x^{-1}$,
\[
  g=e,
\]
contradicción. Así, $\rho_{\mathrm{reg}}(g)$ no fija ningún vector de la base,
y por consiguiente su matriz de permutación tiene diagonal nula. Luego
\[
  \chi_{\mathrm{reg}}(g)=0.
\]

Concluimos que
\[
  \chi_{\mathrm{reg}}(g)=
  \begin{cases}
    |G|, & \text{si } g=e,\\
    0, & \text{si } g\ne e.
  \end{cases}
\]
\hfill$\blacksquare$

\problem{7}{%
Sean $(U,\pi)$ y $(V,\rho)$ dos representaciones de $G$. Demuestre que
$(\Hom_{\CC}(U, V), \tau)$, dada por
\[
  (\tau(g)f)(u) := \rho(g)f(\pi(g^{-1})u),
\]
es, en efecto, una representación.
}

\answer
Para cada $g\in G$ y $f\in\Hom_{\CC}(U,V)$,
\[
  \tau(g)f=\rho(g)\circ f\circ \pi(g^{-1}).
\]
Como $\rho(g)$, $f$ y $\pi(g^{-1})$ son lineales, $\tau(g)f$ también es lineal;
por tanto $\tau(g)f\in\Hom_{\CC}(U,V)$. Además, para $a,b\in\CC$ y
$f_1,f_2\in\Hom_{\CC}(U,V)$,
\[
  \tau(g)(af_1+bf_2)=a\tau(g)f_1+b\tau(g)f_2,
\]
de modo que cada $\tau(g)$ es lineal.

Ahora verificamos la ley de representación. Para todo $f\in\Hom_{\CC}(U,V)$ y
$u\in U$,
\[
  (\tau(e)f)(u)=\rho(e)f(\pi(e^{-1})u)=f(u),
\]
luego $\tau(e)=\Id$.

Si $g,h\in G$, entonces para todo $f\in\Hom_{\CC}(U,V)$ y $u\in U$,
\begin{align*}
  (\tau(g)\tau(h)f)(u)
  &=\rho(g)(\tau(h)f)(\pi(g^{-1})u) \\
  &=\rho(g)\rho(h)f(\pi(h^{-1})\pi(g^{-1})u) \\
  &=\rho(gh)f(\pi((gh)^{-1})u) \\
  &=(\tau(gh)f)(u).
\end{align*}
Por lo tanto,
\[
  \tau(g)\tau(h)=\tau(gh).
\]
En particular, $\tau(g^{-1})$ es la inversa de $\tau(g)$, así que
$\tau(g)\in GL(\Hom_{\CC}(U,V))$ para todo $g\in G$. Concluimos que
$(\Hom_{\CC}(U,V),\tau)$ es una representación de $G$. \hfill$\blacksquare$

\end{document}